\chapter{绪论}

\section{研究背景与意义}
随着航空航天技术与无人机（Unmanned Aerial Vehicle, UAV）平台的飞速发展，遥感影像的数据获取能力呈指数级增长\cite{zhu2017deep}。传统的遥感数据处理范式通常依赖于“端-云”协同架构，即由前端传感器采集海量数据，通过高带宽数据链路回传至地面高性能计算中心进行集中式智能解译。然而，在诸多极具时效性要求的空天任务中（如大范围应急巡航、拒止环境下的自主侦察），受限于通信带宽瓶颈、高昂的传输延迟以及数据链路的脆弱性，传统的“先传后算”模式已难以满足系统的实时闭环控制与自主决策需求\cite{shi2016edge}。在这些场景中，受限于通信带宽瓶颈（星地链路通常仅为数Mbps至数百Mbps）、高昂的传输延迟（卫星链路往返延迟可达数百毫秒至数秒）以及数据链路的脆弱性（易受天气、电磁干扰影响），传统架构已难以胜任\cite{zhang2020edge}。因此，将智能解译算法前置部署至无人机或卫星的板载计算节点，实现机载边缘计算（Airborne Edge Computing），已成为遥感智能化演进的必然趋势\cite{cheng2022comprehensive}。机载边缘计算能够在数据产生的源头直接进行处理，仅将提取的高层语义信息（如目标位置、类别、属性）回传，从而大幅降低对通信带宽的需求，缩短决策时延，提升系统的鲁棒性与自主性。

然而，机载边缘计算面临着严苛的物理与系统约束，这为算法设计带来了巨大的挑战。典型的边缘设备（如NVIDIA Jetson TX2/Orin芯片、Google Edge TPU或FPGA）在功耗（通常在10W-30W之间，与个人计算机数百瓦的功耗相比相差两个数量级）、显存带宽（通常仅为数十GB/s）和算力（INT8算力通常为数十TOPS）上存在严格限制\cite{chen2019deep}。这种约束在飞行平台上更为严苛，因为有限的供电能力、严格的散热要求以及对设备重量和体积的苛刻限制，使得高性能计算设备的部署变得极其困难。例如，典型的无人机平台有效载荷仅为数公斤，留给计算设备的重量预算往往不足500克，这直接限制了芯片的选择和散热方案的设计。与此同时，随着遥感应用向精细化发展，视觉感知任务正不可避免地走向多粒度化（Multi-granularity）。传统的遥感解译往往关注单一任务，如目标检测或土地利用分类。然而，现代复杂应用场景往往需要多粒度信息的协同感知。系统不仅需要获取宏观的实例级位置（目标检测，识别目标的种类、位置、方向、大小等属性），还需要精细的空间像素级轮廓（语义分割，区分建筑物、道路、水体等不同地物的精确边界），甚至需要探究微观的光谱像素级物理属性（高光谱分类，通过光谱特征识别地物的材料成分、健康状态等内在属性）。这种多粒度感知需求使得计算复杂度呈指数级增长，对算力的要求远超单一任务。

这种多粒度复杂感知需求与机载极端算力约束之间的矛盾，构成了当前遥感解译领域的核心技术壁垒。传统的解决方案往往在精度与效率之间进行权衡：要么牺牲精度以换取实时性，采用简化模型降低计算开销；要么保持精度但牺牲实时性，仅能在高性能服务器上运行，无法部署到边缘平台。近年来，以Transformer架构\cite{dosovitskiy2021image}和视觉基础模型（Vision Foundation Models, VFMs，如Segment Anything Model, SAM\cite{kirillov2023segment}）为代表的深度学习算法虽然在感知精度上取得了突破，在ImageNet\cite{russakovsky2014imagenet}、COCO等标准数据集上屡创精度新高，但其动辄数十亿的参数量和数百GFLOPs的计算开销，使其在边缘端的直接部署成为计算上的灾难。以SAM为例，其图像编码器（ViT-H/16）拥有6.36亿个参数，处理一张1024×1024的图像需要约150GFLOPs的计算量，在边缘芯片上的推理时间可达数十秒，完全无法满足实时性要求。

自然界中的生物视觉系统为打破这一算力瓶颈提供了理想的理论启示。生物系统能够在极低能耗下（人脑的功耗仅为20W左右，而GPU的功耗可达数百瓦），通过突触可塑性修剪、频空联合感知以及稀疏选择性注意力等机制，高效地处理海量高维信号\cite{higgins2023towards}。具体而言：（1）大脑发育过程中的突触修剪机制能够自适应地切断不必要的神经连接，实现网络的动态稀疏化，从而降低能耗并提高计算效率；（2）人类视觉系统采用频空联合感知策略，视网膜神经节细胞对不同空间频率具有选择性，能够高效地编码边缘、纹理等关键特征\cite{serre2007robust}；（3）昆虫视觉系统（如蜜蜂复眼）采用光谱选择性机制，仅对特定频段敏感，实现了对复杂光谱信息的稀疏编码。这些生物机制为人工神经网络的设计提供了重要启示。因此，探索仿生视觉感知机理，并将其与深度神经网络的轻量化技术相融合，针对不同初始状态的模型提出定制化的压缩与重构策略，具有重要的理论价值与工程意义。本研究旨在通过构建多维度的仿生轻量化范式，解决严苛算力约束下的多粒度遥感解译难题，为下一代空天一体化自主感知系统的研发提供核心算法支撑。研究成果不仅能够推动遥感智能解译技术的实用化进程，还能够为其他边缘计算场景下的视觉感知任务提供参考和借鉴，具有重要的科学意义和应用前景。

\section{国内外研究现状}
在机载边缘计算场景下，遥感影像的智能解译面临着精度与算力的双重博弈。为了厘清本文研究的切入点，本节将从“遥感多粒度视觉感知算法的演进”、“深度学习轻量化范式”，以及“仿生机理在计算视觉中的应用”三个维度，对国内外现有的研究成果与技术瓶颈进行系统性综述。
\subsection{遥感图像多粒度视觉感知研究现状}
遥感影像的视觉感知任务依据其对地物目标表征精细度的不同，可划分为实例级（目标检测）、空间像素级（语义分割）与光谱像素级（高光谱分类）三个逐步深入的粒度层次。这种多粒度划分反映了遥感应用需求的层次化特征：从粗粒度的"是什么"（目标检测），到中粒度的"在哪里"（语义分割），再到细粒度的"有什么属性"（高光谱分类），不同层次的信息相互补充，共同支撑复杂场景下的智能解译与决策。近年来，随着深度学习技术的突破，各个粒度层次的算法都取得了显著进展，但同时也面临着边缘部署的严峻挑战。

（1）实例级感知：遥感目标检测的发展与瓶颈

目标检测旨在精准定位并识别影像中特定类别的实例对象，输出目标的边界框（Bounding Box）和类别标签。由于遥感影像具有视角高远（通常为几百至数千米的拍摄高度）、目标尺度变化剧烈（同一类别的目标在不同成像距离下占据的像素数量可能相差数十倍）、排列密集（如密集停放的车辆、排列紧密的建筑物）且带有任意方向性（受成像角度影响，目标可能呈现任意旋转角度）等特点，其检测难度远高于自然场景图像。此外，遥感影像中的目标往往缺乏足够的纹理细节，且容易受到阴影、云雾等环境因素的影响，进一步增加了检测难度。

基于卷积神经网络（CNN）的方法： 早期研究多基于通用检测框架进行适应性改进。Ren等人提出的Faster R-CNN\cite{ren2015faster}引入区域建议网络（RPN），实现了端到端的两阶段检测，成为目标检测领域的里程碑工作；Redmon等人提出的YOLO\cite{redmon2016yolo}将检测问题视为单阶段回归任务，大幅提升了检测速度；Liu等人提出的SSD\cite{liu2016ssd}结合了YOLO的速度和Faster R-CNN的精度，通过多尺度特征图实现了对不同尺寸目标的有效检测。这些方法在自然图像数据集上取得了优异性能，但直接应用于遥感影像时效果有限。针对遥感目标的旋转特性，Ding等人提出了经典模型RoI Transformer\cite{ding2019learning}，通过空间变换网络（STN）将水平候选框转换为旋转候选框，显著提升了旋转目标的检测精度；Han等人提出的ReDet\cite{han2021redet}利用旋转等变网络提取具有旋转不变性的特征，设计了面向旋转目标的特征增强模块；Xie等人提出的Oriented R-CNN\cite{xie2021oriented}则设计了定向候选框生成网络，直接生成带有方向信息的建议框，避免了复杂的后处理。

基于Transformer与自注意力机制的方法： 随着Vision Transformer (ViT) 的引入，全局感受野的建模能力得到质的飞跃。Dosovitskiy等人开创性地将Transformer应用于图像识别任务，将图像分割为多个patch并将其视为序列进行建模，实现了与CNN相当甚至更优的性能。Carion等人提出的DETR\cite{carion2020end}打破了传统检测方法依赖非极大值抑制（NMS）的后处理限制，通过Transformer的解码器直接输出检测框，实现了端到端的目标检测，该架构设计简洁优雅，避免了锚框设计和手工特征等繁琐步骤。随后，Zeng等人在遥感领域提出了ARS-DETR\cite{zeng2024ars}，针对遥感目标极端长宽比（如机场跑道、桥梁）的特点，设计了纵横比敏感的注意力机制，有效处理了传统方法难以应对的长条形目标。Transformer方法虽然在精度上表现出色，但其自注意力机制的计算复杂度为O(n²)，对于高分辨率遥感影像而言，计算开销极其庞大。

边缘部署瓶颈： 尽管上述重型检测网络在公开数据集上屡创精度新高，但其特征金字塔（FPN）层级的加深和自注意力矩阵的庞大计算量，导致其参数量动辄上亿，FLOPs高达数百G。以Faster R-CNN-ResNet101为例，其参数量约为6000万，处理一张1024×1024的图像需要约300GFLOPs的计算量。在算力受限的无人机平台上，这些模型极易出现内存溢出（OOM）或推理帧率低下（通常低于1 FPS）。更为严峻的是，当为了适应边缘设备而强行削减网络通道时，高空微小目标在深层网络中的特征响应会迅速衰减甚至弥散，导致严重的漏检。这是因为微小目标在图像中仅占据较少的像素，经过多次下采样后，其特征可能完全消失，而通道剪枝进一步削弱了网络提取细粒度特征的能力。如何在保证检测精度的前提下，实现高效的边缘部署，成为遥感目标检测领域亟待解决的关键问题。

（2）空间像素级感知：遥感语义分割与视觉基础模型

语义分割要求对影像进行像素级的密集预测，为每个像素分配语义类别标签，提取建筑物、道路、水体等空间几何形态，是细粒度解译的核心。与目标检测相比，语义分割提供了更精细的空间信息，能够刻画地物的精确边界和内部结构，在城市规划、土地利用调查、灾害评估等应用中具有不可替代的价值。

全卷积网络与上下文聚合： 自Long等人提出FCN\cite{long2015fully}以来，语义分割进入深度学习时代。FCN将全连接层替换为卷积层，实现了端到端的像素级预测，成为语义分割的奠基性工作。为了解决遥感影像中复杂拓扑结构提取难的问题，DeepLab系列\cite{chen2018deep}利用空洞空间金字塔池化（ASPP）捕获多尺度上下文，通过不同扩张率的空洞卷积提取不同感受野的特征，有效处理了多尺度目标的分割问题；Swin-Unet\cite{cao2022swin}等模型则将Transformer结构引入编码器-解码器架构，利用自注意力机制捕获全局依赖关系，解决了卷积神经网络感受野受限的问题，提升了对大尺度地物（如大型建筑物、农田）的建模能力。Audebert等人提出了"先分割后检测"（Segment-Before-Detect）的思想\cite{audebert2016segmentation}，通过语义分割的结果引导目标检测，提升了检测的准确性和鲁棒性。

视觉基础模型（Foundation Models）的降维打击： 近年来，以Segment Anything Model (SAM)\cite{kirillov2023segment}为代表的零样本（Zero-shot）视觉大模型重塑了分割任务的范式。SAM通过在1100万张图像和10亿个掩码上进行预训练，展现出了惊人的泛化分割能力，能够在无需任何额外标注的情况下，对各类图像进行高质量的分割。这种零样本能力彻底改变了传统的分割任务范式，从"训练-测试"的监督学习模式转向"预训练-提示"的新范式。为了将其引入遥感领域，Cong等人提出了RS-SAM\cite{cong2023rs}，探索了通过提示（Prompt）工程进行遥感影像分割的可行性，设计了针对遥感影像的提示生成策略，验证了SAM在遥感场景中的潜力。

边缘部署瓶颈： SAM庞大的图像编码器使其完全无法在机载芯片上实时运行。SAM-H/16拥有6.36亿个参数，需要约22.6GB的显存来加载模型和中间激活，这远超边缘设备的内存容量（通常为4GB-16GB）。更为致命的是，通用大模型的连续下采样操作（ViT的patch embedding和Transformer的多层下采样）不可避免地平滑了遥感影像中的细粒度边界。在面对极其破碎、纹理高度混淆的复杂形变结构时（如密集的城市建筑群、破碎的地块），SAM常表现出严重的领域偏置（Domain Bias），产生明显的边缘粘连与误判。此外，传统的全量微调（Full Fine-tuning）不仅成本高昂（需要大量标注数据和GPU算力），更易引发灾难性遗忘（Catastrophic Forgetting），即模型在适应新任务时丢失了原有的泛化能力。

（3）光谱像素级感知：高光谱分类的维度灾难

高光谱图像（Hyperspectral Image, HSI）包含了从可见光到近红外波段的数百个连续光谱通道（通常为100-300个波段），每个像素都对应一条完整的光谱曲线，能够捕捉地物微观的物理与化学属性。与RGB图像仅包含3个波段相比，高光谱图像提供了丰富的光谱信息，能够区分在视觉上相似但光谱特征不同的地物，如不同类型的植被、不同材质的人造目标等，在精细农业、矿物勘探、环境监测等领域具有重要应用价值。

从传统降维到深度特征提取： 早期的研究依赖于主成分分析（PCA）或支持向量机（SVM）进行降维与分类。PCA通过线性变换将高维光谱数据投影到低维空间，保留了主要的方差信息；SVM通过核函数将数据映射到高维空间，寻找最优分类边界。然而，这些传统方法难以充分利用高光谱图像的空间信息，且依赖于手工设计的特征。随后，3D-CNN被广泛引入以提取空-谱联合特征，通过3D卷积核同时在空间维度和光谱维度进行卷积，捕获了像素间的空间相关性和波段间的光谱相关性。Li等人对高光谱分类的深度学习方法进行了系统综述\cite{li2019hyperspectral}，总结了从早期光谱分类到空谱联合学习的演进历程。近年来，Hong等人提出的SpectralFormer\cite{hong2021spectral}将Transformer成功应用于高光谱分类，将光谱序列视为一维序列，利用自注意力机制捕获连续光谱波段间的序列依赖关系，在多个基准数据集上取得了优异性能。Zhang等人研究了基于深度学习的卫星图像分类方法\cite{zhang2020satellite}，探索了不同网络架构在遥感场景中的表现。

边缘部署瓶颈： 高光谱数据面临着极其严重的“维度灾难”与Hughes现象。数百个通道的输入张量使得传统深度网络在最前端的计算量呈指数级膨胀。目前主流的高光谱深度网络在机载设备上处理一帧数据往往需要数秒甚至数十秒，完全无法满足实时性要求。如何打破“全波段输入-深层提取”的传统定势，实现原生的计算前置降维，是亟待解决的技术盲区。

% \begin{figure}[htp]
% 	\centering
% 	\includegraphics{picmain}
% 	\caption{图 1.1 名称}
% \end{figure}

\subsection{深度视觉模型的轻量化范式现状}
面对“重型模型算力饥渴”与“边缘设备算力贫瘠”的客观鸿沟，学术界在模型轻量化领域进行了深入探索，逐渐形成了三大主流技术范式：

（1）模型压缩范式（Model Compression）

模型压缩旨在在不显著损失模型精度的前提下，通过去除模型中的冗余信息来降低计算开销。主要包括网络剪枝、知识蒸馏和量化三大技术。

网络剪枝（Pruning）： 分为非结构化剪枝与结构化剪枝。非结构化剪枝剔除权重矩阵中接近零的元素，能够实现很高的压缩比，但生成的稀疏矩阵难以在通用硬件上高效运行，加速效果有限；结构化剪枝则剔除整个卷积核或通道，虽然压缩比略低，但生成的是稠密模型，在CPU、GPU等通用硬件上能够实现显著的加速。Li等人\cite{li2017pruning}提出的滤波器剪枝（Filter Pruning）通过评估通道权重（如L1范数、L2范数）剔除冗余连接，是降低FLOPs最直接的手段。Han等人\cite{han2015deep}提出了深度压缩（Deep Compression）框架，结合剪枝、量化和霍夫曼编码，在ImageNet上实现了35-49倍的压缩。但对于多尺度感知任务，粗暴的静态全局通道剪枝往往会破坏网络对浅层细粒度特征的表征能力，因为微小目标在浅层的特征响应本身就较弱，通道剪枝可能直接切断其特征流路，导致漏检。

知识蒸馏与量化： 知识蒸馏通过大型“教师网络”指导小型“学生网络”的训练。量化则将网络的浮点运算降阶为低精度运算。这些方法虽然有效，但通常作为正交的工程优化手段，并未触及网络结构设计的核心。

（2）参数高效微调范式（Parameter-Efficient Fine-Tuning, PEFT）

随着视觉基础模型（Vision Foundation Models）的兴起，如何高效地将这些预训练大模型适应到下游任务成为研究热点。传统的全量微调（Full Fine-Tuning）需要修改所有模型参数，存储开销巨大（每个任务需要存储一份完整的模型副本），且容易导致灾难性遗忘。PEFT方法通过仅训练少量参数即可实现任务迁移，大幅降低了存储和计算开销。

Adapter与LoRA机制： Houlsby等人提出的Adapter机制\cite{houlsby2019parameter}通过在Transformer块之间插入瓶颈层（Bottleneck Layer，先降维再升维），在冻结主干网络权重的条件下，仅训练极少量参数即可完成下游任务迁移。Adapter的优势在于模块化设计，不同的任务可以共享主干网络，仅替换Adapter模块即可。Hu等人提出的LoRA\cite{hu2022lora}则通过在注意力权重矩阵旁路引入低秩分解矩阵降低了训练开销。具体而言，LoRA假设预训练权重的更新是低秩的，因此用A、B两个低秩矩阵的乘积来近似更新量，可训练参数量仅为原模型的千分之一。LoRA的优势在于推理时可以将低秩矩阵融入原权重，不增加任何计算开销。

遥感适配的局限性： 现有的PEFT方法大多继承自NLP领域，缺乏对视觉底层物理特性的考量。自然语言处理中的任务是离散的（token级别的分类），而视觉任务是连续的（像素级别的预测），且对空间几何、边缘纹理等高频信息极其敏感。直接将常规Adapter应用于遥感图像的大模型微调，难以有效弥补大模型在空间高频细节和细粒度梯度变化上的特征缺失。例如，SAM在遥感影像的细碎边界处往往出现过度平滑，简单的Adapter难以修复这一问题，需要引入针对视觉任务的特殊设计，如空-频联合增强、像素差分等机制。

（3）轻量级原生架构设计范式（Lightweight Native Architecture）

与模型压缩不同，原生轻量化架构从设计之初就考虑了计算效率，通过高效的算子设计、优化的网络结构和合理的模型缩放策略，在保证性能的同时实现低计算开销。这种范式的优势在于无需后处理压缩，模型天生就是高效的。

高效算子设计： 自MobileNet\cite{howard2019searching}系列引入深度可分离卷积（Depthwise Separable Convolution）以来，轻量化网络迎来了爆发。深度可分离卷积将标准卷积分解为深度卷积（Depthwise Convolution，每个卷积核只处理一个输入通道）和逐点卷积（Pointwise Convolution，1×1卷积用于通道融合），计算量从标准卷积的O(K²C²HW)降低到O(K²CHW + C²HW)，当卷积核K较大时，加速效果显著。随后，ShuffleNet\cite{ma2019shuffle}利用通道重排（Channel Shuffle）促进组卷积的信息交互，避免了组间信息孤立的问题；EfficientNet\cite{tan2019efficient}通过神经架构搜索（NAS）系统性地探索了网络深度、宽度、分辨率的缩放策略，实现了精度-效率的最优平衡。这些算子大幅降低了特征编码的乘加操作数（MACs），使得在移动设备上实时运行深度网络成为可能。

多模态/高维数据的适用性： 现有的轻量化算子多针对RGB三通道图像设计，通过减少通道间的交互来降低计算量。然而，当面对包含数百通道的高光谱数据时，简单的深度可分离卷积依然会引发巨大的通道内计算量。高光谱数据的特殊性在于：（1）波段数量远超RGB图像，通道间相关性复杂，简单的深度卷积难以有效建模；（2）不同波段的重要性差异巨大，存在大量冗余波段；（3）光谱曲线是连续的，需要建模波段间的序列依赖关系。必须结合数据的物理先验，重新设计动态稀疏的原生架构。例如，可以借鉴生物视觉的光谱选择性机制，在网络最前端动态筛选重要波段，实现计算前置降维；也可以设计空-谱联合卷积，利用通道间的相关性进行信息聚合，降低通道内的计算开销。

\subsection{仿生机理在计算视觉中的应用}
将自然界生物系统“节能与高效”的机制抽象为数学模型并引入计算视觉，是当前人工智能前沿的重要交叉方向。

神经突触可塑性与动态稀疏： 在大脑发育过程中，突触修剪（Synaptic Pruning）机制自适应地切断微弱或不必要的突触，从而维持神经网络的高效与低能耗。研究表明，人类大脑在青春期会经历大规模的突触修剪，删除约40%的突触连接，这个过程不是随机的，而是遵循"用进废退"的原则：经常使用的突触会被强化，而很少使用的突触会被删除。利用这种机制指导模型进行动态稀疏化训练，使网络能够在推理时自动关闭无关特征通道，实现"按需计算"。具体实现方法包括：（1）基于重要性评分的剪枝，如基于梯度、权重幅值、一阶泰勒展开等指标评估通道的重要性；（2）基于正则化的稀疏训练，通过L1正则、组稀疏正则等促使网络学习稀疏权重；（3）基于奖励机制的强化学习，通过Agent决定每个神经元是否激活，实现自适应的稀疏计算。这些方法都受到生物突触可塑性的启发，旨在在保证性能的前提下，最大化网络的稀疏性。

人类视觉系统（HVS）的频空联合感知： 人眼视网膜神经节细胞对空间频率呈现出高度的选择性，倾向于利用低频信号构建全局结构，利用高频信号敏锐捕捉边缘纹理。具体而言，视网膜上的细胞分为M细胞（大细胞，Magno）和P细胞（小细胞，Parvo）：M细胞对低空间频率和高时间频率敏感，负责运动检测和粗糙形状识别；P细胞对高空间频率和颜色敏感，负责精细细节和纹理识别。这种"双通路"设计使得人类视觉能够在不同分辨率下高效地处理信息。引入傅里叶变换模拟该机制，能够打破传统卷积操作过度关注局部空域像素的局限。具体方法包括：（1）频域特征增强，在傅里叶域对高频分量进行加权，增强边缘信息；（2）多频带特征融合，将图像分解为不同频带，分别提取特征后融合；（3）可学习频域滤波器，设计可学习的频域卷积核，自适应地捕获不同频带的特征。这些方法都借鉴了人类视觉的频-空协同机制，能够在提升边缘检测精度的同时，降低计算开销。

昆虫视觉的光谱选择性与多尺度聚合： 蜜蜂复眼中的光感受器仅对特定频段敏感，且不同的小眼负责不同尺度的视野特征聚合。这种“稀疏光谱编码”机制为解决高光谱数据的计算前置降维提供了绝佳的仿生学模板。

\section{本文面临的主要挑战与科学问题}
在机载边缘计算的极限算力与内存带宽约束下，实现多粒度高精度感知，本文致力于解决以下三个核心科学问题。这些问题反映了当前遥感解译领域在理论方法和技术实现层面的核心瓶颈，也是推动边缘智能感知技术突破的关键所在。

科学问题一：在进行高压缩比的网络结构剪枝时，如何打破“参数量下降”与“微小目标多尺度特征保真度”之间的相互制约？
重型检测模型在被强行“瘦身”后，往往会不可逆地丢失高空微小目标的特征响应。如何在剔除冗余通道的同时，动态重构特征路由路径，避免“一剪即损”，是实现高效目标搜寻的关键难题。

科学问题二：在严禁修改视觉大模型主体权重的前提下，如何以极低的算力开销补偿并增强网络对复杂高频边缘信息的感知能力？

通用大模型（如SAM）的连续下采样操作不可避免地平滑了遥感影像中的细粒度边界。SAM的ViT编码器采用patch embedding（如16×16的patch），相当于进行了16倍下采样，这直接抹除了小于16×16像素的细节信息。对于遥感影像中经常出现的细碎地物（如密集建筑群、小块农田、细长道路），这种下采样导致严重的边缘模糊。然而，直接微调SAM的全部参数面临巨大挑战：（1）计算资源需求巨大，SAM-H拥有6.36亿参数，全量微调需要数十块高端GPU；（2）存储开销高昂，每个下游任务需要存储一份完整模型；（3）容易引发灾难性遗忘，过度的参数调整可能导致模型失去原有的零样本泛化能力。因此，必须设计一种轻量化的旁路适配机制，在不修改主体权重的前提下，利用频域-空域协同补充细粒度特征。具体而言，需要解决以下子问题：如何设计高效的频域增强模块来恢复边缘信息？如何通过像素差分等操作显式捕捉高频纹理？如何将这些增强特征以最小的计算开销注入到冻结的主干网络中？这是实现基础大模型低成本机载泛化的核心挑战。

科学问题三：面对包含数百维度光谱信号的严重计算冗余，如何跳出传统的后处理降维框架，从架构原生层面实现空谱联合特征提取的极致轻量化？
传统高光谱模型多采用将全量数据输入深层网络再进行特征选择的逻辑。如何建立一种“计算前置”的动态稀疏筛选机制，从源头掐断冗余光谱信息的计算流，是突破极端轻量化高光谱分类（KB级算力）的理论瓶颈。

\section{本文主要研究内容与贡献}
针对上述科学问题，本文以“机载算力受限场景下的多粒度视觉感知”为核心研究主线，基于“模型压缩”、“大模型微调”、“极简架构设计”三种轻量化范式，开展了递进式的算法创新研究。本文的主要研究内容与贡献如下：

（1）面向实例级感知：提出了基于突触分形剪枝的轻量化遥感目标检测方法（LiteFM-RSDet）

主要贡献： 针对现有检测模型参数冗余与微小目标检测精度下降的矛盾，设计了突触分形剪枝机制（Synaptic Fractal Pruning, SFP），创新性地将分形数学中的自相似递归结构引入网络剪枝。SFP利用分形门控函数递归评估特征通道的信息熵，模拟大脑突触的"用进废退"机制，自适应地切断冗余特征流路。与传统的全局统一剪枝不同，SFP能够根据特征图的空间位置和尺度差异进行差异化剪枝，最大限度地保留微小目标的关键特征。同时，提出了**多分辨率流形路由（Multi-resolution Manifold Routing, MMR）**策略，通过拉普拉斯曲率感知，将不同尺度的目标特征路由至轻量级专家模块。MMR借鉴了流形学习的思想，在特征空间中自适应地划分不同尺度的流形，为每个尺度分配专用的轻量专家，避免了单一专家难以同时处理多尺度目标的问题。在DOTA、HRSC2016等遥感目标检测数据集上的实验表明，该模型将参数量从原始的100+M大幅压缩至28.3M，推理速度提升4倍以上，同时保持了与原模型相当的检测精度，突破了轻量化模型对微小目标多尺度感知能力衰退的瓶颈。

（2）面向空间像素级感知：提出了结合频空协同的大模型参数高效微调方法（FourierPDC-SAM）

主要贡献： 针对SAM边缘分割模糊且微调成本高昂的问题，提出了像素差分卷积（Pixel Difference Convolution, PDC）模块，通过计算像素邻域的梯度差分显式增强高频纹理捕捉能力。PDC受到人类视觉皮层中边缘检测细胞的启发，通过差分操作直接建模像素间的变化关系，突破了传统卷积对局部相似性的依赖。引入了基于傅里叶变换的频域增强机制（Fourier Enhancement, FE），分离低频全局结构与高频干扰噪声。FE将图像变换到频域，针对性地增强边缘对应的高频分量，同时抑制低频背景噪声，模拟了人类视觉系统中P细胞对高频信号的选择性关注。通过轻量化适配器将频空特征注入冻结的SAM主干，整个适配器的参数量仅为SAM主干的0.1%，但显著提升了分割精度。在xBD、LoveDA等遥感语义分割数据集上的实验表明，FourierPDC-SAM在xBD数据集上实现了IoU 3.3%的显著提升（从71.2%提升至74.5%），推理速度提升3倍，实现了视觉大模型向遥感细粒度分割的极低成本迁移。更重要的是，该方法保持了SAM的零样本泛化能力，能够处理训练未见过的目标类别。

（3）面向光谱像素级感知：提出了模拟蜜蜂视觉机制的原生极简高光谱分类架构（BioLiteNet）

主要贡献： 针对高光谱数据的维度灾难，创新性地提出了蜂感选择器（BeeSense Selector, BSS），模拟蜜蜂寻找花蜜时的稀疏注意机制，在网络最前端直接过滤掉90\%以上的冗余波段，实现计算前置降维。BSS借鉴了蜜蜂复眼对特定光谱波段敏感的特性，通过可学习的注意力权重动态选择最重要的几个波段，从源头掐断了冗余计算流。联合设计了多尺度仿射卷积（Affine Scale Convolution, ASC）进行空谱解耦运算，ASC通过仿射变换建模光谱维度的非线性变化，同时利用多尺度卷积核捕获空间结构信息。整个网络的参数量仅为传统3D-CNN的1\%，但分类精度相当。在Indian Pines、Pavia University、Salinas Scene等多个高光谱基准数据集上的极限测试表明，该架构成功将计算复杂度极致压缩至~90K FLOPs（KB级别），推理时间从传统的数百毫秒降低至毫秒级，在保持95\%以上精度的同时，确立了边缘高光谱感知的新极限。BioLiteNet证明了通过原生架构设计而非后处理压缩，可以实现真正的极致轻量化。

\section{论文组织结构}
本文全文共分为六章，具体组织结构安排如下：

第一章 绪论： 阐述课题背景与系统约束，系统梳理国内外在多粒度感知、模型轻量化及视觉大模型微调领域的现状，分析现有方法的局限性，明确本文的核心科学问题与主要贡献，为后续研究奠定基础。

第二章 多粒度感知理论与仿生轻量化方法基础： 详细阐述检测、分割、分类任务的数学表征与优化目标，定义模型复杂度的评价体系（包括参数量、FLOPs、内存占用、推理延迟等指标），梳理依托的仿生学理论基础（突触可塑性、人类视觉频空感知、昆虫光谱选择性），介绍相关的数学工具和理论框架。

第三章 实例级感知（模型压缩范式）： 论述基于SFP与MMR的轻量化遥感检测方法（LiteFM-RSDet）的设计原理与实验论证。首先详细介绍突触分形剪枝和多分辨率流形路由的技术细节，然后在DOTA、HRSC2016等数据集上进行充分实验，并与现有方法进行对比分析，验证该方法的有效性和优越性。

第四章 空间像素级感知（大模型微调范式）： 剖析结合PDC与FE的SAM大模型参数高效微调方法（FourierPDC-SAM）。阐述像素差分卷积和频域增强的设计动机和实现方式，介绍轻量化适配器的网络结构。在xBD、LoveDA等遥感语义分割数据集上进行实验，评估方法的精度提升和计算开销，并对消融实验结果进行分析。

第五章 光谱像素级感知（原生轻量设计范式）： 详细展开极简高光谱分类架构（BioLiteNet）及其蜂感选择机制，验证其在极限算力约束下的有效性。介绍蜂感选择器和多尺度仿射卷积的设计原理，在Indian Pines、Pavia University、Salinas Scene等多个高光谱基准数据集上进行广泛实验，对比分析与其他方法的性能差异，论证原生架构设计的优势。

第六章 总结与展望： 对全文研究成果进行系统性总结，归纳三大轻量化范式（模型压缩、大模型微调、原生轻量化）的理论价值和技术贡献，分析当前方法的局限性和适用场景，并对基于多模态大模型及神经形态计算的边缘感知发展方向进行展望，为后续研究提供参考。


